\section{Optimal Transport 视角}

轨迹弯曲的根本原因在于 $\mathbf{x}_0$ 和 $\mathbf{x}_1$ 的\textbf{配对方式}。本节从 Optimal Transport（最优传输）的视角来理解这一问题。

\subsection{随机配对 vs.\ 最优配对}

在标准 Flow Matching 训练中，$\mathbf{x}_0$ 和 $\mathbf{x}_1$ 是\textbf{独立}采样的：
$$\mathbf{x}_0 \sim p_0, \quad \mathbf{x}_1 \sim p_1, \quad \text{两者独立}$$

这意味着一个噪声样本 $\mathbf{x}_0$ 可能被``分配''到任意一个数据样本 $\mathbf{x}_1$，没有空间结构上的偏好。

\begin{importantbox}{最优传输（Optimal Transport）的基本思想}
最优传输问题关注的是：如何找到从分布 $p_0$ 到分布 $p_1$ 的最``经济''的传输方案？具体地，对于 $L_2$ 代价（又称 Wasserstein-2 距离），最优传输方案最小化：
$$W_2^2(p_0, p_1) = \inf_{\pi \in \Pi(p_0, p_1)} \int \|\mathbf{x}_0 - \mathbf{x}_1\|^2 \, d\pi(\mathbf{x}_0, \mathbf{x}_1)$$
其中 $\Pi(p_0, p_1)$ 是所有耦合（联合分布）的集合，其边际分别为 $p_0$ 和 $p_1$。

\medskip
\textbf{直觉}：最优传输倾向于将``近处''的点配对，避免长距离搬运，从而最小化总搬运代价。
\end{importantbox}

\subsection{OT 配对对轨迹直度的影响}

如果我们用最优传输方案来配对 $\mathbf{x}_0$ 和 $\mathbf{x}_1$，而不是随机配对，那么：

\begin{itemize}
    \item 配对更``合理''：附近的噪声点倾向于映射到附近的数据点
    \item 不同配对之间的条件轨迹交叉更少
    \item 边际向量场的冲突更少，边际轨迹自然更直
\end{itemize}

\begin{warningbox}{OT 配对的计算代价}
尽管 OT 配对在理论上很有吸引力，但精确计算 Optimal Transport 方案的复杂度很高。对于 $n$ 个样本，经典算法的复杂度为 $O(n^3)$，这在大规模训练中不太实际。因此，实际中通常采用近似方法，例如在 mini-batch 内计算 OT（Mini-batch OT），或者使用 Reflow 作为替代方案——Reflow 可以看作是一种``隐式的'' OT 配对优化。
\end{warningbox}

\subsection{Reflow 与 OT 的关系}

Reflow 和 Optimal Transport 之间有深刻的联系：

\begin{itemize}
    \item Reflow 通过迭代微调，隐式地优化了 $\mathbf{x}_0 \to \mathbf{x}_1$ 的配对方案
    \item 每一轮 Reflow 都使配对更接近 OT 方案——因为模型学到的映射倾向于走``最短路径''
    \item 理论上，无限次 Reflow 迭代的极限就是 OT 映射（在某些正则性条件下）
\end{itemize}

\begin{knowledgebox}{从传输代价看轨迹弯曲}
轨迹弯曲的本质是\textbf{传输代价过高}。当远处的点被配对在一起时，连接它们的直线会穿越中间区域，与其他配对的直线交叉。交叉导致边际向量场不一致（同一点处有多个方向的速度），网络只能学习期望（平均）速度，结果就是弯曲的轨迹。OT 配对通过最小化传输代价来减少交叉，Reflow 则通过迭代隐式实现类似效果。
\end{knowledgebox}

\subsection{本章小结}

Optimal Transport 提供了理解轨迹弯曲问题的理论框架。随机配对导致高传输代价和轨迹交叉，OT 配对可以缓解这一问题但计算代价高。Reflow 提供了一种实用的替代方案，通过迭代微调隐式优化配对。

